\scanpage{350}
Sous ces conditions, pour tout \(U\in\operatorname{Ob}\unsure{\mathscr X|X|\mathcal X}\),
\[
 \inserted{C^*(U)=}\Gamma(\deleted{\unsure{U}}\operatorname{Hom}(U,C^*))=
 \Gamma(U,C_U^*)
\]
est \deleted{un}\inserted{\(k(U)\)-}complexe \(\uline{k(U)}=\)\deleted{\(\Gamma(U,k_U)\)}-\uline{augmenté} muni des \unsure{mêmes types} de structure, ce qui implique que son \deleted{\(\unsure{H^*(U,C_U^*)}\)}\inserted{\(H^*=H^*(C^*(U))\)} est muni d'une structure \unsure{évidente} d'algèbre \inserted{\uline{graduée}} associative \unsure{unit. anticommutative}. Si \(C^*\) est une \uline{résolution} de \(k\), \unsure{on en déduit qu'on a} un hom
\[
 H^*(C^*(U))\longrightarrow H^*(U,k_U)\tag*{(\(\unsure{*|1}\))}
\]
\margin{\uline{\unsure{la prem.|le prou.|En fait}}}
% slanted
\unsure{et on vérifie que c'est} un hom d'anneaux avec \unsure{unité}. Ce dernier \unsure{fait} implique à son tour une \uline{\unsure{preuve}} d'\unsure{unicité}.

\uline{Ex.1} Cas où \(C^*\) a une structure d'Alg. diff.

\uline{Ex.2} Nous nous \unsure{trouverons au cas où on a|sommes en un cas où on a} un \(k\)-Module. \unsure{Soit} une \unsure{surjection} dans \(\Phi\) définie par une section \(T\) de \(\Phi\). On suppose que \(\Phi/k.T=\Psi\) est \unsure{\(k\)-plat}. On sait alors que le complexe
\[
 \begin{aligned}
 \cdots\longrightarrow C^0\longrightarrow C^1\longrightarrow\cdots
   \longrightarrow C^n\longrightarrow C^{n+1}\longrightarrow\cdots\\
 \shortparallel\qquad\quad\shortparallel\qquad\qquad\qquad
   \shortparallel\qquad\quad\shortparallel\\[-2pt]
 0\longrightarrow\Phi\xrightarrow{T\wedge}\bigwedge^2\Phi
 \xrightarrow{T\wedge}\cdots\longrightarrow\bigwedge^{n+1}\Phi
 \xrightarrow{T\wedge}\bigwedge^{n+2}\Phi\longrightarrow\cdots
 \end{aligned}
\]
est une \uline{résolution} de \(k\), avec
\[
 Z^i\simeq\bigwedge^i\Psi\qquad(i\geq0)
\]
(ce qui exprime que \unsure{l'on a une} \unsure{suite} exacte \unsure{canonique|canoniquement}
\[
 0\longrightarrow\bigwedge^{i-1}\Psi\longrightarrow
 \bigwedge^i\Phi\longrightarrow\bigwedge^i\Psi\longrightarrow0\qquad).
\]
% archive sheet number: 351
